\documentclass[10pt,a4paper]{article}
\usepackage[margin=21mm]{geometry}
\usepackage{fontspec}
\setmainfont{lmroman10-regular.otf}[BoldFont=lmroman10-bold.otf,ItalicFont=lmroman10-italic.otf,BoldItalicFont=lmroman10-bolditalic.otf]
\usepackage{amsmath,amssymb,booktabs,tabularx,array,graphicx,caption,fancyhdr,xcolor}
\usepackage[authoryear,round]{natbib}
\usepackage[unicode,colorlinks=true,linkcolor=black,citecolor=black,urlcolor=blue]{hyperref}
\usepackage{xurl,microtype}
\DeclareMathSizes{10}{10}{7}{7}
\setlength{\parindent}{0pt}\setlength{\parskip}{4pt}
\setlength{\emergencystretch}{3em}
\renewcommand{\arraystretch}{1.08}
\captionsetup{font=footnotesize,labelfont=bf}
\setcounter{topnumber}{4}\renewcommand{\topfraction}{.95}\renewcommand{\textfraction}{.06}\renewcommand{\floatpagefraction}{.75}
\setlength{\textfloatsep}{10pt}\setlength{\intextsep}{10pt}
\pagestyle{fancy}\fancyhf{}\fancyhead[L]{\small Hakan Zeki Gülmez}\fancyhead[R]{\small SBA 7(a) lending | Working paper}\fancyfoot[C]{\thepage}\setlength{\headheight}{13pt}
\definecolor{blue}{HTML}{0072B2}
\newcommand{\losslabel}{EAD = GrossApproval; full-disbursement proxy, CCF = 100\%. LGD is the gross charge-off proxy on the same approval denominator. SBA/lender amounts use assumption-based pro-rata guarantee allocation.}
\newcommand{\scenario}{Scenario projections on the fixed 2006 portfolio; not out-of-sample validation.}
\newcommand{\tbl}[2]{\begin{table}[htbp]\centering\input{generated/#1.tex}\caption{#2}\end{table}}
\newcommand{\datafile}[1]{\href{https://github.com/hakangulmez/credit-risk-sba}{\nolinkurl{#1}}}
\input{generated/numbers.tex}
\input{generated/editorial_numbers.tex}

\hypersetup{pdftitle={Local economic conditions and public-lender loss sharing in SBA 7(a) lending},pdfauthor={Hakan Zeki Gülmez},pdfsubject={Working paper: temporal calibration and conditional loss-sharing scenarios}}
\makeatletter
\setlength{\@fptop}{0pt}
\setlength{\@fpsep}{14pt}
\setlength{\@fpbot}{0pt plus 1fil}
\makeatother

\newcommand{\inlinetbl}[2]{\par\medskip\noindent\begin{minipage}{\linewidth}\centering\input{generated/#1.tex}\captionof{table}{#2}\end{minipage}\par\medskip}

\begin{document}
\begin{center}
{\LARGE\bfseries Local economic conditions and public--lender loss sharing in SBA 7(a) lending}\par\vspace{6pt}
Hakan Zeki Gülmez\par
M.Sc. Management \& Technology (Economics \& Econometrics), Technical University of Munich\par
8 October 2026 \quad $\cdot$ \quad Working paper\par\small \url{https://github.com/hakangulmez/credit-risk-sba}
\end{center}
\begin{abstract}
We study the temporal calibration of credit-risk models and conditional public--lender loss sharing in approved and disbursed SBA 7(a) loans. Official administrative records define the outcome as recorded charge-off within 36 months of first disbursement. Term-free Firth logistic models have modest out-of-time discrimination: crisis AUC is \AucCrisis. Mean predicted crisis risk is \PdCrisis\%, compared with \RateCrisis\% recorded; later-cohort risk is \PdLater\%, compared with \RateLater\%. A separate macro-conditioned model associates a one-percentage-point increase in window-end unemployment relative to its starting value with \Endpoint\ probability percentage points on a fixed reference portfolio. Baseline and adverse expected-loss proxies are \ElBase\ and \ElAdverse\ million dollars under full-disbursement, gross-severity and pro-rata guarantee assumptions. The evidence highlights calibration failure and the dependence of scenario accounting on probability and loss assumptions. Future realized macro paths, unverified feature vintages, revised data and a post-results design amendment limit interpretation. The analysis estimates neither causal guarantee effects nor actual fiscal costs.
\end{abstract}
\noindent\textbf{Keywords:} SBA lending; credit-risk calibration; temporal validation; Firth logit; local economic conditions; public guarantees.\par
\noindent\textbf{Research status:} Working paper based on the saved results of the final estimation round. This presentation revision retains the empirical estimates.

\section{Introduction}
Publicly guaranteed lending connects private credit allocation to public exposure. A lender evaluates applicants and originates loans, while a public guarantee allocates part of the covered exposure to an agency. Understanding this arrangement requires more than sorting borrowers by a risk score. Probability estimates enter loss calculations in levels; a model that preserves some ordering but misstates the frequency of adverse outcomes can yield misleading portfolio amounts. Local economic conditions may also change during the life of a loan, making the information available at approval different from the conditions under which losses occur.

This paper asks: how well do reconstructed approval-descriptor models discriminate and calibrate across disbursement cohorts, and how do conditional macroeconomic scenarios change assumption-based public--lender loss allocation? We study official SBA 7(a) administrative records, rather than treat a selected resolved-loan teaching extract as the complete lending population. The outcome is recorded charge-off within 36 months of first disbursement. A complete window is required for all retained statuses, including unresolved records. The estimand is therefore an administrative fixed-horizon outcome among approved and disbursed loans; it is not borrower distress, delinquency or credit access among rejected applicants.

The empirical design separates two layers. Layer 1 uses reported approval/application descriptors and state indicators, excluding reported maturity from primary specifications because its approval-time vintage cannot be verified. Its evaluation compares an early estimation period with crisis and later cohorts. Layer 2 adds unemployment and house-price information spanning the subsequent observation window. Realized paths support a retrospective conditional cohort check; stipulated paths support a fixed-portfolio scenario exercise. Neither use supplies a real-time macro forecast or a causal intervention.

Four findings organize the analysis. First, out-of-time discrimination is modest. Second, the probability scale does not transport satisfactorily: predicted mean risk is below the recorded crisis rate and above the recorded later-cohort rate. Third, macroeconomic conditional associations depend on how a path is summarized, with endpoint and peak unemployment changes reported together. Fourth, expected-loss and public--lender allocations change under fixed scenarios, but their interpretation remains conditional on exposure, gross severity and guarantee assumptions. Calibration error can affect levels, differences and ratios; a precise conditional interval does not validate those assumptions.

The contribution is an auditable empirical illustration at the intersection of temporal model evaluation and guaranteed-loan scenario accounting. It combines the official status universe, complete-horizon sample accounting, explicit information sets, archived estimation failures, calibration diagnostics and paired scenario uncertainty. It does not introduce a new estimator, establish that machine learning solves credit prediction, or estimate a welfare effect of SBA guarantees. Its value is to expose where an apparently operational credit-risk calculation depends on unverified timing and probability assumptions.

\textbf{Design history.} The final estimation protocol was frozen after first-round results had been seen and before any final-round model was estimated. Deviations are logged in DECISIONS.md. Ordinary-logit fits and calibration results had already been inspected. Training-frequency pooling, Firth estimation and the expanded Layer 2 chronology were then documented before any final-round refit. The amendment is not a confirmatory replication; joint changes cannot identify which isolated change improved or worsened performance. This manuscript uses saved aggregate outputs only and adds no empirical estimation or specification search.

\section{Related literature and economic framework}
\subsection{Public guarantees and the distinction between behaviour and accounting}
\citet{bachas2021} exploit notches in the SBA guarantee schedule to study credit supply. Their design uses policy variation to address lender behaviour. The present paper addresses a different question: how recorded charge-off probabilities and assumed loss allocation vary across cohorts and scenarios in an observed portfolio. Guarantee share in our descriptor model is not an identified treatment, and the pro-rata allocation formula is not evidence that a guarantee changed underwriting, monitoring or lending volume. Policy design and allocation accounting need separate evidence.

A guarantee can change contractual exposure without fully revealing who bears the eventual economic loss. Actual allocation may depend on recoveries, servicing, eligibility, guarantee purchase and subsequent transactions. The available fields establish approved amounts and recorded gross charge-off information, not those cash-flow histories. We therefore use an explicit allocation convention and interpret its outputs as proxies. The SBA share is the scenario-loss-weighted mean of loan guarantee shares; it can change when scenario probabilities reweight loans even though contractual shares are held fixed.

The official-source choice also matters for the risk population. \citet{li} present a loan dataset and teaching framework, while the earlier resolved-loan extract discussed in the retained source audit conditions on terminal status. Retaining mature unresolved records in the official source avoids that particular selection rule. It does not remove selection into approval, establish that unresolved loans perform well, or verify every recorded predictor's original vintage. The distinction concerns population measurement rather than a claim of universal data completeness.

\subsection{Discrimination, calibration and estimation under separation}
\citet{vancalster2019} distinguish risk ranking from agreement between predicted and observed risks and discuss calibration deterioration across settings. Their application domain is clinical prediction; the statistical distinction applies to any probability model, including the credit-risk exercise here. AUC summarizes ordering, while portfolio loss calculations depend on the probability scale. Average agreement also leaves room for errors within risk groups. This motivates reporting AUC, mean predicted/recorded rates, jointly estimated calibration diagnostics and saved training-score groups side by side.

Firth's bias reduction modifies the likelihood score \citep{firth1993}; the finiteness and shrinkage properties of the Jeffreys-prior penalty are studied by \citet{kosmidis2021}. These results support an estimator when ordinary-logit coefficients are not well defined under separation. They do not establish out-of-time calibration. \citet{puhr2017} show that coefficient bias reduction and accurate probabilities are separate concerns, including a tendency toward equiprobability in their rare-event settings. That finding motivates caution; it does not establish that the Firth penalty causes the cohort errors observed here. Differences in conditions, composition, measurement or programme practices remain observationally entangled.

\subsection{Scenario conditioning and economic interpretation}
The distinction between realized information and a stipulated path is central. A model conditioned on unemployment and housing over a loan's future window can describe associations or translate assumptions into scenario probabilities. It cannot be evaluated as a forecast made at disbursement when those future observations were unknown. Likewise, a fixed historical crisis path is an input to a scenario exercise rather than an estimated probability distribution of future shocks.

The \citet{bcbs2018} principles emphasize clear objectives, methodology and documentation in stress-testing frameworks. We use that guidance as a conceptual reference for transparent scenario construction, not as evidence of regulatory compliance. There is no balance-sheet capital model, endogenous lending response or joint economic-path simulation here. Keeping the reference portfolio fixed provides a controlled accounting comparison, while limiting generalization to portfolios with other composition, exposures or guarantee terms.

\section{Data, target and sample}
\subsection{Official sources and administrative measurement}
The official SBA 7(a) FOIA snapshot is dated 30 June 2026. Fiscal-year resource names refer to approval, whereas estimation roles use first disbursement. The dictionary includes payoff and charge-off dates, approval and guaranteed amounts, status, jobs reported on application, business age/type, industry, project state and reported maturity. Current lender assignment and borrower identifiers are not predictors. Source terms designate U.S. Government Works; no borrower-level data are distributed. BLS unemployment is monthly seasonally adjusted, in percent; FHFA HPI is the quarterly, non-seasonally-adjusted all-transactions state index. FRED identifiers are references, not acquired observations.

FHFA observations and previously retrieved BLS observations are unchanged. Only missing unemployment observations needed through the end of the later validation horizon were added in the final estimation round. This creates a mixed retrieval vintage of revised agency histories, not historical release information. Retrieval dates and source hashes are in the frozen source manifest. This product uses FHFA data but is neither endorsed nor certified by FHFA.

For disbursement date $d_i$, define $w_i=d_i+36$ calendar months, with month-end clamping. We retain a complete observation window through the snapshot and define
\[Y_i=\mathbf{1}\{\text{valid ChargeOffDate}_i\in[d_i,w_i]\}.\]
The outcome is \emph{recorded charge-off within 36 months}. Mature PIF and EXEMPT observations without such a charge-off receive zero. EXEMPT means unresolved status withheld under the statutory exemption; it is not certified performing or free of delinquency. Late charge-offs are outside the horizon rather than removed from the administrative population. Amount anomalies affect severity checks separately from PD eligibility.

\tbl{waterfall}{Sequential source-population eligibility accounting; overlapping source flags are not substituted for this waterfall. Counts are saved research aggregates.}
\tbl{samples_paper}{Layer-specific disbursement roles, events/non-events and complete macro coverage. The original heldout-check role is post-amendment retrospective validation, with the disclosed cohort reuse. The fixed scenario portfolio is included in expanded Layer 2 estimation. Event counts use the administrative fixed-horizon definition.}

\subsection{Feature timing and comparability}
Approval/application amounts and jobs are treated as reported descriptors, with jobs reflecting lender estimates rather than realized employment. We use log approval, guarantee share, log-one-plus jobs and categorical processing, rate type, sector, franchise presence, business age/type, loan type and project state. Historic vintages of several descriptors remain unverified. The initial interest rate is excluded from shared specifications for comparability because missingness is extensive; we do not estimate a historical interest-rate association. Reported Term is excluded from every primary model, regardless of the audit verdict.

\tbl{rate_missingness}{InitialInterestRate missingness; excluded from all shared specifications for comparability, not silently imputed as an observed interest rate.}

Numeric medians, missing flags and standardization use only the corresponding estimation sample. Within each nonstate categorical variable, levels whose training count share is strictly below the frozen threshold are pooled to Other. No outcome counts choose the maps; states are never pooled with divisions or each other. Unseen levels use a fitted Other category where supported; otherwise no unsupported dummy is introduced. Native tree categories use their documented missing fallback. Constant or exactly redundant columns are omitted without outcome selection. Maps, reference categories, rank/separation checks and application counts remain available in the frozen outputs.

\section{Chronology, models and inference}
\subsection{Different information sets and temporal roles}
Layer 1 estimates on 1991--2002, with benchmark tuning in the first estimation round on 2003H1 and benchmark calibration on 2003H2. ElasticNet and LightGBM hyperparameters and selected boosting rounds were inherited unchanged in the final estimation round: no new search or early stopping occurred. The crisis test is 2007--2009; the additional cohort is 2011--2013. Calibration diagnostics for Firth do not modify its raw probabilities.

Layer 2 estimates on 1991--2009. The LightGBM robustness uses the corresponding frozen Layer 1 parameters/rounds from the first estimation round, and its probability calibrator is fitted on 2010 only. Validation is 2013--2014: \textbf{Retrospective conditional validation using realized macro paths.} Estimation labels mature by end-2012; calibration labels extend through end-2013 and overlap the calendar period of the 2013 validation cohort. That cohort was already included in first estimation round evaluation. This is post-amendment validation, not a previously untouched confirmatory test. No held-out result changes the specification or threshold.

\subsection{Firth and ordinary-logit comparison}
For $p_i=\Lambda(x_i'\beta)$, $\eta_i=x_i'\beta$ and $\Lambda(z)=1/(1+e^{-z})$, the primary estimate maximizes
\begin{align*}
\ell^*(\beta;w)&=\sum_iw_i\{Y_i\eta_i-\log(1+e^{\eta_i})\}+\tfrac12\log\det I_w(\beta),\\
I_w(\beta)&=X'\operatorname{diag}\{w_ip_i(1-p_i)\}X.
\end{align*}
The point fit uses unit weights, without oversampling, class balancing or intercept adjustment. Adjusted-score iteration uses the frozen convergence/conditioning criteria. Layer 1 contains state effects among origination categories. Layer 2 adds macro inputs and a linear disbursement-year trend centred on 2000, rather than unidentified future-year dummy coefficients. Corresponding ordinary-logit comparators use identical rows and reduced columns; invalid fits do not provide valid coefficients or predictions.

The final estimation round converged for each Firth point fit. Layer 1 ordinary MLE has an outcome-pure pooled category even though numerical optimization reports success; Layer 2 reaches the fixed iteration limit, which alone does not prove separation. The with-Term ElasticNet benchmark also fails at its frozen limit. These failures remain visible; no alternative favourable specification replaces them. Complete fit diagnostics and implementation validation are linked in the appendix.

\tbl{convergence_paper}{Saved Firth and same-design ordinary-logit status; numerical optimizer success is not a valid-estimate guarantee.}

\subsection{Calibration diagnostics and their exact definition}
For a frozen probability vector $\widehat p_i$, the saved diagnostic regression jointly estimates $\alpha$ and $b$ in
\[\Pr(Y_i=1\mid\widehat p_i)=\Lambda\{\alpha+b\,\operatorname{logit}(\widehat p_i)\}.\]
The reference pair is $(0,1)$. The reported intercept is from this joint regression: it is not an intercept-only calibration-in-the-large estimate with the slope fixed at one. Its sign alone therefore does not determine whether mean risk is over- or underpredicted. We assess that direction directly from saved mean probabilities and recorded cohort rates. These diagnostics evaluate predictions without changing the primary Firth probabilities. Neither ideal joint parameters nor mean agreement alone would establish calibration in every subgroup.

\subsection{Two uncertainty procedures}
Stored evaluation intervals use state-cluster resampling conditional on a fitted model; these include discrimination, calibration intercept/slope and group-rate bands. They do not include coefficient estimation uncertainty. Training-score decile boundaries are frozen from the fit sample, not chosen from evaluation scores. We reuse saved edges, unequal counts, empty bins and undefined values without new binning or binomial intervals.

Firth training-refit inference instead uses positive state-weighted multiplier attempts. Each scheduled draw assigns independent Exponential(1) weights to states, copies them to their loans and normalizes within the estimation sample to its original loan count. Every state remains present; underlying weights are shared across specifications and paired scenarios. The weighted likelihood and Firth penalty are re-estimated in the final estimation round, not at reporting time. Linear percentile bounds use only successful fully defined draws and report their effective denominators. No failed attempt is redrawn.

\tbl{bootstrap_paper}{Scheduled training-state multiplier refits: all attempts retained, no redrawing. Evaluation bands separately use 999 state-cluster draws conditional on fitted models.}

State independence and conditioning on successful fits are assumptions. Preprocessing, portfolio, LGD, exposures, allocation shares and paths remain fixed. Source-vintage, calibration-model and scenario-construction uncertainty are excluded. Calibration error can alter nonlinear PD levels, differences and ratios, so inference does not validate the probability scale.

\section{Discrimination and calibration results}
\tbl{performance}{All valid Term-free primary/benchmark models in the selected evaluation cohorts. AP denotes average precision, not a different risk population. Saved 95\% intervals condition on each fitted model; no model is selected for reporting because it scores highest.}
\tbl{calibration}{Mean predicted/recorded rates and diagnostic calibration intercept/slope, with stored 95\% intervals. Ideal diagnostic values are zero/one; Firth predictions stay raw. Layer 2 is retrospective conditioning on future realized paths.}

Layer 1 crisis AUC is \AucCrisis\ and additional-cohort AUC \AucLater. The ranking evidence is above chance but modest; it is not strong or stable transportability. Mean Firth risk is \PdCrisis\% against \RateCrisis\% recorded in the crisis and \PdLater\% against \RateLater\% later. Conditions, composition and programme or underwriting changes are possible explanations, not identified mechanisms. The macro-tree calibration failure in the first estimation round remains documented in the retained first-round report; the revised chronology does not remove that evidence.

For Layer 2, raw Firth predicts \LtwoFirth\% against \LtwoRate\% recorded. Calibrated LightGBM predicts \LtwoTree\%, closer in average, but its diagnostic intercept \TreeIntercept\ and slope \TreeSlope\ exclude their respective ideal values. We therefore describe average agreement as positive descriptive evidence, not calibration solved. Figure~\ref{fig:cal} and the saved group tables show the distinction.

In the 2007--2009 crisis cohort, the lowest training-score group has \CrisisGroupOnePred\% predicted versus \CrisisGroupOneRecorded\% recorded. The two highest groups have nearly equal recorded rates (\CrisisGroupNineRecorded\% and \CrisisGroupTenRecorded\%) despite mean predictions of \CrisisGroupNinePred\% and \CrisisGroupTenPred\%. Descriptively, the error is more than a common level shift: the score separates recorded rates little at the top, consistent with the stored diagnostic slope of \CrisisDiagnosticSlope. No difference test or mechanism is inferred.

For calibrated Layer 2 LightGBM in 2013--2014, predicted and recorded rates are close descriptively in training-score groups 1--8. Larger deviations appear in the two highest groups: \TreeGroupNinePred\% versus \TreeGroupNineRecorded\%, and \TreeGroupTenPred\% versus \TreeGroupTenRecorded\%. The latter contains only \TreeGroupTenN\ loans, with a stored recorded-rate band of [\TreeGroupTenLower, \TreeGroupTenUpper]\%. Average agreement coexists with overprediction in the highest-risk groups; the wide top-group band limits precision. These comparisons use saved groups and bands, without new intervals or tests.

\begin{figure}[htbp]\centering\includegraphics[width=\linewidth]{figures/calibration_groups.pdf}\caption{Saved training-score decile diagnostics. Original bin numbers/counts/edges and empty bins are retained in the accompanying tables. Stored rate bands are state-cluster bootstrap intervals conditional on fitted models, not independent-loan binomial CIs. Layer 2 uses future realized paths.}\label{fig:cal}\end{figure}

\section{Conditional macro associations}
For origin month $m(d_i)$ and containing quarter $q(d_i)$, define
\begin{align*}
u_0&=u_{m(d_i)},&\Delta u&=u_{m(w_i)}-u_{m(d_i)},\\
h_0&=100\log(H_{q(d_i)}/H_{q(d_i)-4}),&\Delta h&=100\log(H_{q(d_i)+12}/H_{q(d_i)}).
\end{align*}
Unemployment changes are in percentage points. Quarterly HPI is not interpolated; containing-period alignment is an approximation to exact-day windows. The separately estimated sensitivity replaces only $\Delta u$ by $\text{peak\_du}=\max_{m=m(d_i),\ldots,m(w_i)}(u_m-u_0)$, including origin and endpoint.

We evaluate the finite change
\[\frac{100}{N_R}\sum_{i\in R}\{\Lambda(\eta_i+\beta_{\Delta u}^{\rm native})-\Lambda(\eta_i)\},\]
on the fixed observed 2006 reference portfolio, holding every other covariate constant. The native coefficient removes train-standardization scaling. This is a +1 pp contrast, not an infinitesimal derivative or an intervention on unemployment. The endpoint result \Endpoint\ and peak sensitivity \Peak\ are presented together. A rise that later reverses can vanish from endpoint change; this motivates the predeclared sensitivity rather than post-results selection.
\tbl{macro_paper}{Fixed observed portfolio, conditional +1 pp finite-change associations. Probability percentage points; training-refit state-multiplier intervals conditional on fixed inputs. Peak is a separately estimated sensitivity, not a replacement.}

\section{Scenario projections and loss accounting}
The fixed 2006 scenario loans also enter expanded Layer 2 estimation. Baseline uses each state's 2004--2006 mean unemployment level, zero unemployment change, mean annual HPI log growth for $h_0$, and three times that growth for $\Delta h$. Adverse uses fixed January 2007--January 2010 unemployment and 2007Q1--2010Q1 housing paths, without choosing a worst date. Full scenario vectors remain in the frozen path table. These paths are conditions; support checks do not establish economic plausibility.

For segment severity $L_{s(i)}$, approved exposure $A_i$ and guarantee share $g_i$,
\[EL_i=p_i L_{s(i)} A_i,\qquad g_i=\frac{\text{SBAGuaranteedApproval}_i}{A_i},\]
\[EL^{SBA}=\sum_i g_i EL_i,\qquad EL^{lender}=\sum_i(1-g_i)EL_i.\]
\losslabel\ We estimate segment LGD as the ratio of training gross charge-off sums to approval sums, using loan type and fixed approval-size cells with the retained sparse-cell pooling: insufficient cells fall back to loan-type severity, then the overall training severity, rather than an outcome-selected new segmentation. Primary severity is uncapped; recovery amounts and amortization are absent. Full utilization is not measured exposure and does not imply that the combined EL estimate is necessarily conservative. Neither allocation represents observed payouts, measured public spending or lender behaviour.

\tbl{stress}{\scenario\ PD and loss rates are percentages, changes in PD are percentage points, and amounts are USD millions. Difference bounds come from paired draws, never subtraction of marginal interval endpoints. Assumption-based loss proxies; see Section 7.}

Baseline/adverse loss proxies are \ElBase\ and \ElAdverse\ million dollars; the arithmetic ratio is \ElRatio. The stored paired change is \ElDiff\ million. SBA allocation shares are \ShareBase\% and \ShareAdverse\%; lender shares are their complements. Different scenario PDs reweight loans, so these aggregate shares need not be fixed. No share or ratio interval is invented.

\section{Sensitivities and support}
\inlinetbl{sensitivities_paper}{Stored adverse-path sensitivities. With-Term is timing-unverified; peak substitutes for endpoint change. Capped LGD caps each training severity amount at its own approval before unchanged pooling; downturn LGD uses only the original designated recession-exposed cohorts. \scenario\ Assumption-based loss proxies; see Section 7.}
Downturn severity uses 1991, 2000 and 2001 training cohorts only, without enlarging the window. The old-extract utilization/CCF sensitivity remains unavailable because licence, cutoff and population admissibility were not established. Urban/rural classification is absent, and missing historical rates preclude comparable rate associations. These unavailable items are not replaced by new calculations. Full with-Term results remain accessible as qualified sensitivities rather than headline validation evidence.

\inlinetbl{support}{State-specific versus pooled-global univariate training-range checks on primary inputs. Loan and approval exposures are union counts, not sums of repeated feature flags. \scenario\ Exposure is an approval-dollar count, not an estimated loss.}
The state-specific adverse flagged share is \SupportShare\%; global flags are zero. This can reflect broad pooled ranges even where state-specific values extrapolate. Neither univariate check establishes joint support or shows whether mixed baseline/adverse paths are economically plausible. Complete feature-level flags are linked to the frozen support output.

\section{Economic discussion}
\subsection{A probability model can fail in different directions across cohorts}
The central economic result is the failure of probability calibration across periods. The crisis and later cohorts exhibit errors in opposite directions. This matters whenever a score is multiplied by an exposure or used to compare proposed lending policies. A modest AUC does not supply a probability adjustment, and finite coefficients do not establish operational suitability. The evidence supports reporting temporal calibration as a first-order model property rather than judging the exercise by its best ranking metric.

We cannot identify which mechanism drives the shifts. Local conditions, borrower composition, underwriting, programme changes, administrative recording and unobserved borrower balance sheets may all matter. A comparison of Layer 1 and Layer 2 does not isolate a macro contribution: their estimation windows and information sets differ. Nor does average agreement for calibrated LightGBM prove that recalibration solves the problem. Its saved joint diagnostic and risk-group evidence remain imperfect, and calibration uses a distinct cohort with the disclosed calendar overlap. These findings motivate monitoring and properly timed evaluation, rather than an immediate deployment recommendation.

\subsection{Macro paths describe conditional exposure to local conditions}
A positive endpoint-unemployment association is consistent with worse repayment conditions in weaker local economies, but the present design does not separate that channel from correlated housing, industry, borrower selection or omitted financial conditions. State effects remove fixed differences represented by those indicators; they do not make subsequent unemployment changes exogenous. The reported finite contrast holds other regressors fixed and does not reproduce the full economic adjustment that an intervention or recession would generate.

Endpoint and peak definitions answer different descriptive questions. An unemployment increase that later reverses may be invisible in the endpoint difference yet present in the peak statistic. Reporting both prevents a single summary from being treated as the entire path. Housing inputs also summarize a containing-quarter approximation rather than exact daily economic conditions. The associations are useful for understanding how a chosen conditional model translates paths into probabilities, while their economic interpretation remains limited by those summaries and information assumptions.

\subsection{Public--lender loss sharing depends on weights and loss measurement}
The pro-rata SBA amount is a sum of guaranteed shares times model-based loan losses. It illustrates how a contractual allocation convention interacts with portfolio composition and scenario-specific probabilities. Similar aggregate shares in the two scenarios do not demonstrate a behavioural invariance, and their difference need not be zero. Loans with different shares receive different risk weights when the path changes. An intervention that raised guarantee shares could also change lending and borrower selection, which this fixed-portfolio accounting does not model.

The dollar amounts should be read together with calibration evidence. A systematic probability distortion can move both scenario levels and the paired contrast; the logit's nonlinearity provides no automatic cancellation. The interval for the contrast captures successful coefficient refits under fixed portfolio and loss assumptions. It does not bound uncertainty about actual utilization, net recoveries, revised predictor histories or fiscal payments. Expected losses also do not establish the distribution of aggregate losses or its tail. The accounting exercise consequently cannot rank the social desirability of alternative guarantee policies.

\section{Limitations and future work}
Measurement is administrative rather than economic delinquency. Mature EXEMPT loans are unresolved, not certified performing. Only approved/disbursed loans enter the universe. Jobs are lender-reported application estimates; several descriptor vintages and reported Term remain unverified. Initial rates are missing over much of the estimation period. Current data revisions and containing-period macro alignment cannot reconstruct origination information sets.

Layer 2 sees future realized paths in cohort checks, with post-results amendments, a previously inspected cohort and calibration-calendar overlap. Raw Firth calibration fails, and average agreement for a calibrated tree cannot establish group calibration. Gross approval, full utilization and gross severity omit recovery/amortization information; pro-rata guarantees are assumed. Intervals assume independent state clusters and use only the refits that converged (198 of 199 for the primary scenario model). Fixed loss assumptions and univariate support checks further limit their interpretation.

\textbf{Discussion-only extensions.} Rolling-origin recalibration would separate recalibration from refitting, use only outcomes whose windows had matured at each origin, and predeclare methods/dates after disclosing inspected cohorts. Two-block Shapley accounting would compare unemployment and housing blocks including initial levels/changes, averaging both orders across four fixed coefficient/scenario combinations; it would not be causal attribution and would require hybrid-path support checks. Competing risks could distinguish charge-off and payoff, using existing payoff dates after assessing date meaning and follow-up; PIF need not be early repayment, and Term vintage remains unresolved. Additional local-sector/lender information and nonlinear/hierarchical PD or severity models would require separate timing/value assessments, not automatic market expansion.

Layer 3 is deferred to separately authorized design review. A policy design would need verified facts, prior-outcome-exposure disclosure and credible comparison. DML, doubly robust DiD or causal forests cannot supply absent identification; a documented rejected design is acceptable. No such analysis is performed here.

\section{Conclusion}
The approval-descriptor scorecard retains modest ranking but fails probability calibration across cohorts. Conditional macro associations depend on path definition and future information. The loss-sharing exercise converts those probabilities and frozen severity/exposure assumptions into scenario accounting rather than fiscal measurement. Taken together, the findings make temporal calibration central to interpreting guaranteed-loan risk and show how conditional loss allocation inherits probability and measurement assumptions. Further progress requires better timed validation and exposure evidence, rather than stronger claims from the current dollar projections.

\begin{thebibliography}{99}
\bibitem[Bachas et al.(2021)]{bachas2021} Bachas, N., Kim, O. and Yannelis, C. (2021). Loan guarantees and credit supply. \emph{Journal of Financial Economics}, 139(3), 872--894. \url{https://doi.org/10.1016/j.jfineco.2020.08.008}.
\bibitem[BCBS(2018)]{bcbs2018} Basel Committee on Banking Supervision (2018). \emph{Stress testing principles}. Bank for International Settlements, October. \url{https://www.bis.org/bcbs/publ/d450.htm}.
\bibitem[BLS(2026)]{bls} Bureau of Labor Statistics (2026). Local Area Unemployment Statistics: official API and metadata. \url{https://www.bls.gov/lau/data.htm}. Frozen retrieval details: source manifests; revised agency history.
\bibitem[FHFA(2026)]{fhfa} Federal Housing Finance Agency (2026). All-Transactions House Price Index, states, not seasonally adjusted. \url{https://www.fhfa.gov/data/hpi/datasets}. Required attribution notice appears above.
\bibitem[Firth(1993)]{firth1993} Firth, D. (1993). Bias reduction of maximum likelihood estimates. \emph{Biometrika}, 80(1), 27--38. \url{https://doi.org/10.1093/biomet/80.1.27}.
\bibitem[Kosmidis and Firth(2021)]{kosmidis2021} Kosmidis, I. and Firth, D. (2021). Jeffreys-prior penalty, finiteness and shrinkage in binomial-response generalized linear models. \emph{Biometrika}, 108(1), 71--82. \url{https://doi.org/10.1093/biomet/asaa052}.
\bibitem[Li et al.(2018)]{li} Li, M., Mickel, A. and Taylor, S. (2018). Should this loan be approved or denied? A large dataset with class assignment guidelines. \emph{Journal of Statistics Education}. \url{https://doi.org/10.1080/10691898.2018.1434342}. Bibliographic/source rights verified in retained DATA.md; dataset mirror admissibility remains unresolved.
\bibitem[Puhr et al.(2017)]{puhr2017} Puhr, R., Heinze, G., Nold, M., Lusa, L. and Geroldinger, A. (2017). Firth's logistic regression with rare events: accurate effect estimates and predictions? \emph{Statistics in Medicine}, 36(14), 2302--2317. \url{https://doi.org/10.1002/sim.7273}.
\bibitem[SBA(2026)]{sba} U.S. Small Business Administration (2026). 7(a) \& 504 FOIA dataset, snapshot 30 June 2026. \url{https://data.sba.gov/dataset/7a-504-foia}. Original source-file hashes retained; 7(a) only.
\bibitem[Van Calster et al.(2019)]{vancalster2019} Van Calster, B., McLernon, D.J., van Smeden, M., Wynants, L. and Steyerberg, E.W. (2019). Calibration: the Achilles heel of predictive analytics. \emph{BMC Medicine}, 17, 230. \url{https://doi.org/10.1186/s12916-019-1466-7}.
\end{thebibliography}

\appendix
\section{Saved estimation and Term diagnostics}
\tbl{rank}{Outcome-independent rank reduction and retained state effects. A reference state is represented by the intercept. Complete column maps and coefficients remain in saved outputs.}
\tbl{coefficients}{Selected coefficients from the primary/sensitivity specifications. Numeric columns use training standardization; associations are not causal. The original full coefficient table is linked below. Stored training-refit intervals hold other uncertainty components fixed.}
\tbl{term}{All applicable eligible date-valid charge-offs/PIF, including charge-offs after the label window. Matching is reported Term minus elapsed months within three months; elapsed months are actual days divided by the average calendar-month length $365.25/12$ days fixed in the protocol. Invalid reported Term does not establish its vintage.}
Verdict: \emph{Timing unverifiable.} Low charge-off date matching (\TermMatch\%) cannot verify approval-time measurement; supporting metadata or an original vintage is absent. PIF coincidence also does not prove updating. Every primary specification remains Term-free. Missing/invalid counts and approval-year/loan-type distributions are available in the complete frozen audit.

\section{Calibration-group detail and reproducibility}
The following tables preserve training-score cutoffs, unequal group counts and empty/unavailable values. Rate bands use stored conditional state-cluster evaluation inference, not independent-loan binomial CIs. Intercepts/slopes in the main table include stored intervals; fitted-cohort diagnostics are not substituted for these evaluation checks.
\tbl{bins_1}{Layer 1 raw Firth crisis cohort. Training-score edges in probability percent; recorded-rate bands condition on the fitted model.}
\tbl{bins_2}{Layer 1 raw Firth additional cohort. Same saved training-score grouping convention.}
\tbl{bins_3}{Layer 2 raw Firth validation cohort; future realized macro paths, not real-time prediction.}
\tbl{bins_4}{Layer 2 calibrated LightGBM validation cohort. Original training-score groups; average agreement does not establish risk-group calibration.}

\textbf{Complete audit links.} \href{https://github.com/hakangulmez/credit-risk-sba/blob/main/results/g2-bis-2026-10-08/term_audit_summary.csv}{Term distributions}, \href{https://github.com/hakangulmez/credit-risk-sba/blob/main/results/g2-bis-2026-10-08/term_audit_status_counts.csv}{status counts}, \href{https://github.com/hakangulmez/credit-risk-sba/blob/main/results/g2-bis-2026-10-08/fit_status.csv}{fit diagnostics}, \href{https://github.com/hakangulmez/credit-risk-sba/blob/main/results/g2-bis-2026-10-08/firth_coefficients.csv}{full coefficients}, and \href{https://github.com/hakangulmez/credit-risk-sba/blob/main/results/g2-bis-2026-10-08/scenario_support_by_feature.csv}{feature support flags}.\par\textbf{Local evidence links.} Full aggregates: \nolinkurl{results/g2-bis-2026-10-08/}. Audit: \nolinkurl{term_audit_summary.csv}, \nolinkurl{term_audit_status_counts.csv}, \nolinkurl{term_audit_evidence.json}. Methods: \nolinkurl{docs/PROTOCOL_v2_2026-10-08.md}, \nolinkurl{config/g2_bis_2026-10-08.yaml}, \nolinkurl{bootstrap_summary.json}, \nolinkurl{fit_status.csv}, \nolinkurl{firth_coefficients.csv}, \nolinkurl{scenario_support_by_feature.csv}. The public results notebook is available at \url{https://github.com/hakangulmez/credit-risk-sba/blob/main/notebooks/research_walkthrough.ipynb}. Exact unique selectors/hashes and display conversions: \nolinkurl{CLAIMS.md} and \nolinkurl{results/g5-2026-10-08/claims_registry.json}. Saved values and bands are unchanged; a results-only build works without private loans/models/predictions. The historical full test run is retained evidence, distinct from the current reporting checks.
\end{document}
